\chapter{Current Results and Discussion}

\section{Complexity Summary}

The current backends expose the same public operations but have different theoretical costs.
Table~\ref{tab:complexity-summary} summarizes the documented bounds for heap-native operations.

\begin{table}[H]
    \centering
    \begin{tabular}{llllll}
        \toprule
        Backend        & Insert      & Decrease    & Remove          & Peek   & Extract         \\
        \midrule
        Binary heap    & $O(\log n)$ & $O(\log n)$ & $O(\log n)$     & $O(1)$ & $O(\log n)$     \\
        Fibonacci heap & am. $O(1)$  & am. $O(1)$  & am. $O(\log n)$ & $O(1)$ & am. $O(\log n)$ \\
        Kaplan heap    & am. $O(1)$  & am. $O(1)$  & am. $O(\log n)$ & $O(1)$ & am. $O(\log n)$ \\
        \bottomrule
    \end{tabular}
    \caption{Complexity summary for heap-native public operations.}
    \label{tab:complexity-summary}
\end{table}

\texttt{contains} is intentionally excluded from the table.
It is currently linear in all backends and should not be used as a performance-critical operation.

\section{Validation Status}

The repository currently supports the following checks:

\begin{itemize}
    \item normal C test suite;
    \item debug checks with internal invariants enabled;
    \item AddressSanitizer and UndefinedBehaviorSanitizer;
    \item leak-check target when supported by the environment;
    \item optimized test and smoke benchmark run;
    \item Doxygen documentation build;
    \item whitespace check in CI.
\end{itemize}

This is a strong validation baseline for an experimental C project.
It does not prove absence of bugs, but it covers many classes of regression: API edge cases, stale handles, structural invariants, memory safety, undefined behavior, and benchmark smoke consistency.

\section{Engineering Status}

Several pieces make the repository more robust than a minimal algorithm demonstration:

\begin{itemize}
    \item centralized overflow checks for allocation size arithmetic;
    \item generational handles with wrong-heap and stale-handle rejection;
    \item backend-specific invariant checkers;
    \item node pools for metadata reuse;
    \item fallback behavior for temporary consolidation allocation failures;
    \item deterministic randomized tests;
    \item TSV benchmark output and comparison tooling;
    \item hardened DIMACS parser.
\end{itemize}

These features support the experimental goal.
When a backend changes, the repository has tools to evaluate correctness, memory behavior, and performance-sensitive scenarios through the same public API.

\section{Expected Performance Trade-Offs}

The binary heap should be expected to perform well on many practical workloads because it uses contiguous storage and simple loops.
Even though its decrease-key is $O(\log n)$, the constant factors and cache locality may be favorable.

The Fibonacci and Kaplan heaps should be most interesting on workloads with many decrease-key operations.
Their theoretical advantage is amortized decrease-key behavior.
However, they use pointer-heavy structures, node metadata, and consolidation tables.
Therefore, their practical performance depends heavily on workload and machine characteristics.

The Dijkstra benchmark is useful because it exercises a real mixture of insert-handle, decrease-key, and extract-min operations.
The heap-only benchmark is useful because it isolates specific operation patterns.
Both are needed.

\section{Smoke Tests Versus Experimental Results}

The repository includes a tiny DIMACS graph for smoke testing.
This is appropriate for CI because it is fast and deterministic.
It is not a serious performance dataset.

Large road-network datasets can be used locally.
For stronger experimental claims, the project should define a fixed benchmark suite, record machine and compiler information, run multiple repetitions, and store summarized results in the report or in a separate tracked experiment note.
